\documentclass[../../main.tex]{subfiles}
\begin{document}

\begin{center}
  {\fontsize{48pt}{58pt}\selectfont \textbf{2002}\normalsize\textbf{年度}}
\end{center}

\begin{table}[htb]
  \centering
  \caption{出題テーマと難易度}
  \begin{tabular}{|c|p{9cm}|c|}
    \hline
    \textbf{問題番号} & \textbf{テーマ} & \textbf{難易度} \\
    \hline
    第1問 & 和の極限を条件とする数列 & \\
    \hline
    第2問 & 円に内接する三角形の辺長不等式 & \\
    \hline
    第3問 & 整数係数四次方程式の根 & \\
    \hline
    第4問 & 逆双曲線関数の導関数と螺旋の長さ & \\
    \hline
    第5問 & 三次関数の3交点の位置 & \\
    \hline
    第6問 & 複素数の回転で作る点列の周期 & \\
    \hline
  \end{tabular}
\end{table}

\begin{table}[htb]
  \centering
  \caption{解答}
  \renewcommand{\arraystretch}{1.4}
  \begin{tabular}{|c|c|p{9cm}|}
    \hline
    \textbf{問題番号} & \textbf{小問} & \textbf{答え} \\
    \hline
    第1問 & -- & $a_1=1$，$a_2=-1$，$n\ge3$では$a_n=\dfrac{1}{(n-1)(n-2)}$ \\
    \hline
    第2問 & (1) & 辺長の平方和が$8$より大きければ鋭角三角形（証明問題） \\
    \cline{2-3}
    & (2) & 平方和は$9$以下．等号は正三角形 \\
    \hline
    第3問 & -- & $(a,b,c)=(\pm2,2,\pm2),(\pm1,0,\pm1),(\pm3,4,\pm3)$（複号同順） \\
    \hline
    第4問 & (1) & $f'(x)=\dfrac{1}{\sqrt{1+x^2}}$ \\
    \cline{2-3}
    & (2) & $\dfrac{\pi\sqrt{1+\pi^2}+\log(\pi+\sqrt{1+\pi^2})}{2}$ \\
    \hline
    第5問 & -- & $a^2>b$，3交点の$x$座標はすべて$(-a-2\sqrt{a^2-b},-a+2\sqrt{a^2-b})$内（証明問題） \\
    \hline
    第6問 & -- & $z_n=z_0$となる$n\ge1$が存在する条件は$\theta$が有理数（証明問題） \\
    \hline
  \end{tabular}
\end{table}

\paragraph{試験形態}
6問．試験時間は要確認．

\end{document}
